% Source: Lecture17_18.pdf, pages 1-3 (Lecture 17).
% Title on the source: "Lecture 17 - Renormalization Group, continued".
\lecture{17}
\subsection{A finite-temperature fixed point: the spin-1 Ising model}

Here we will discuss an example of the RG with an RG flow having a finite-temperature fixed point (RGFP).
Consider a spin-1 Ising model with FM spin--spin coupling $J>0$ on a 2D (square) lattice. The minimal
spin--spin interaction is
\begin{equation}
  \Ham_{\min} = -J\sum_{e_{ij}}s_is_j,
\end{equation}
where each $s_i$ can take on the values $s_i = -1,0,1$. More generally, $\Ham$ can also contain the n.n.\
term $-J_{sq}\sum_{e_{ij}}s_i^2s_j^2$ and also the single-site term $-J_{ss}\sum_is_i^2$. So with all of
these, one has
\begin{equation}
  -\beta\Ham = K\sum_{e_{ij}}s_is_j + K_{sq}\sum_{e_{ij}}s_i^2s_j^2 + K_{ss}\sum_is_i^2 .
\end{equation}
The additional terms are both invariant under the same global $\mathbb{Z}_2$ invariance as the first term,
so these are irrelevant operators. By the same kind of real-space block spin construction as before, one
can obtain a set of RG equations. Define the Boltzmann variables
\begin{equation}
  x = e^{-(K+K_{sq}+K_{ss}/2)},\qquad y = e^{-(K+K_{sq}+K_{ss})},\qquad z = e^{-2K}.
\end{equation}
The $z$ variable is the same as before, and the $x$ and $y$ variables arise from the two additional
couplings. Note that there is no external magnetic field; $H = 0$. As before, the RG transformation
maintains this property. The RG equations with a blocking parameter $b = 2$ are
\begin{equation}
  x' = \left[\frac{x(1+y+z)}{1+x^2+z^2}\right]^2,\qquad
  y' = \left[\frac{y^2+2x^2}{1+x^2+z^2}\right]^2,\qquad
  z' = \left[\frac{2z+x^2}{1+x^2+z^2}\right]^2
\end{equation}
(details in J.~Yeomans, \emph{Statistical Mechanics of Phase Transitions}, with notation changes).

We observe that if we start with $x = 0$, the RG transformation $\Rightarrow x' = 0\ \forall\,y,z$, so we
simplify the analysis by restricting to the plane $x = 0$ in this 3-dimensional space of variables. Then the
RG transformation becomes
\begin{equation}
  z' = \frac{4z^2}{(1+z^2)^2},\qquad y' = \frac{y^4}{(1+z^2)^2}.
\end{equation}
Since we consider $J>0$, so $K\ge 0$, the variable $z$ is in the range $0\le z\le 1$, and $y$ is finite.
N.B.: $z = 0\Rightarrow z' = 0$ and $z = 1\Rightarrow z' = 1$ (independent of $y$), and
$y = 0\Rightarrow y' = 0$, independent of $z$, so (a) any point on the left, right, or bottom side stays on
that side; (b) the set of RGFPs includes $(z_*,y_*) = (0,0)$ and $(1,0)$. With $z = 0$, we get $y' = y^4$;
the physical solution is $y = 1$, so another RGFP is $(0,1)$.

Next, set $y = 0$ (so $y' = 0$) and solve for $z$ at an RGFP, $z' = z$:
\begin{equation}
  z = \frac{4z^2}{(1+z^2)^2}.
\end{equation}
One solution is $z = 0$, but we already included that, so take $z\ne 0$ and divide by $z$
$\Rightarrow 1 = 4z/(1+z^2)^2$, $1+2z^2+z^4 = 4z$, $z^4+2z^2-4z+1 = 0$. This equation factors:
\begin{equation}
  (z-1)(z^3+z^2+3z-1) = 0 .
\end{equation}
The solution $z = 1$, $y = 0$: $(1,0)$ is already listed, so we just have to solve the cubic equation
$z^3+z^2+3z-1 = 0$. There is a unique physical solution, $z = 0.295598\ldots$, so this gives
$(z_*,y_*) = (0.2956,0)$. (The other two solutions of the cubic are complex: $z = -0.6478\pm 1.7214\ii$.)
We will round this off to $(z_*,y_*) = (0.30,0)$.

Next, set $z = 1 = z'$, so this will be an RGFP; solve for $y$: $y' = y^4/(1+z^2)^2 = y^4/4$,
so $y' = y\Rightarrow y = y^4/4$, $y^4-4y = 0$, $y(y^3-4) = 0$; the solution with $y = 0$, namely
$(z_*,y_*) = (1,0)$, has already been listed, so take $y\ne 0$. The cubic equation $y^3 = 4$ has a unique
physical solution, $y_* = 4^{1/3} = 2^{2/3} = 1.58740$. (The other 2 solutions are complex:
$y = 2^{2/3}e^{2\pi\ii/3}$ and $y = 2^{2/3}e^{-2\pi\ii/3}$.) So another RGFP is
$(z_*,y_*) = (1,1.5874)\simeq(1,1.59)$. Now setting $z = 0.2956$ yields $z' = z$, so we can find another RGFP
by solving the equation for $y' = y^4/(1+z^2)^2$ with this value of $z$. This gives
\begin{equation}
  (z_*,y_*) = (0.2956,1.0574)\simeq(0.30,1.06).
\end{equation}

With the same methods that we have discussed before, we can determine the RG flow for all $z$ and $y$, as
shown in the diagram.

\begin{diagram}
\begin{tikzpicture}[scale=1.0,>=Stealth]
  \pgfmathsetmacro{\sx}{5} \pgfmathsetmacro{\sy}{2.4}
  \draw (0,0) rectangle (\sx,2.2*\sy);
  % edges: z = 0, y = 0 and z = 1 are invariant lines
  \draw[->] (0,{0.5*\sy}) -- (0,{0.25*\sy}); \draw[->] (0,{1.35*\sy}) -- (0,{1.6*\sy});
  \draw[->] (0.297*\sx,0) -- (0.15*\sx,0); \draw[->] (0.297*\sx,0) -- (0.65*\sx,0);
  \draw[->] (\sx,{1.2*\sy}) -- (\sx,{0.7*\sy}); \draw[->] (\sx,{1.7*\sy}) -- (\sx,{2.0*\sy});
  % the line z = 0.2956 is invariant; the flow on it leaves the red RGFP
  \draw[thick] (0.2956*\sx,0) -- (0.2956*\sx,2.2*\sy);
  \draw[->,thick] (0.2956*\sx,{0.75*\sy}) -- (0.2956*\sx,{0.5*\sy});
  \draw[->,thick] (0.2956*\sx,{1.4*\sy}) -- (0.2956*\sx,{1.8*\sy});
  % flow lines leave the red RGFP (repulsive along both eigenvectors); computed from the map,
  % the thick ones are the separatrices to (0,1) and (1,1.59)
  \draw[->] plot[smooth] coordinates {(1.478,2.538) (1.365,2.513) (1.270,2.494) (1.168,2.477) (1.073,2.463) (0.978,2.450) (0.874,2.440) (0.779,2.433) (0.675,2.430) (0.579,2.436) (0.487,2.462) (0.407,2.527) (0.361,2.611) (0.329,2.710) (0.309,2.804) (0.294,2.899) (0.281,3.002) (0.272,3.098) (0.264,3.193) (0.257,3.297) (0.251,3.393) (0.245,3.497) (0.241,3.593) (0.236,3.688) (0.232,3.792)};
  \draw[] plot[smooth] coordinates {(0.232,3.792) (0.229,3.888) (0.226,3.992) (0.223,4.088) (0.220,4.184) (0.218,4.288) (0.216,4.384) (0.213,4.480) (0.211,4.584) (0.209,4.680) (0.208,4.784) (0.206,4.880) (0.204,4.976) (0.203,5.080) (0.201,5.176) (0.200,5.280)};
  \draw[->] plot[smooth] coordinates {(1.478,2.538) (1.377,2.515) (1.291,2.498) (1.204,2.483) (1.117,2.472) (1.029,2.467) (0.941,2.473) (0.859,2.502) (0.791,2.557) (0.742,2.630) (0.704,2.718) (0.678,2.802) (0.657,2.887) (0.640,2.974) (0.626,3.061) (0.614,3.148) (0.603,3.235) (0.594,3.323) (0.585,3.410) (0.578,3.498) (0.570,3.594) (0.564,3.681) (0.558,3.769) (0.553,3.857) (0.548,3.945)};
  \draw[] plot[smooth] coordinates {(0.548,3.945) (0.544,4.033) (0.539,4.121) (0.535,4.209) (0.531,4.297) (0.528,4.384) (0.524,4.480) (0.521,4.568) (0.518,4.656) (0.515,4.744) (0.512,4.832) (0.510,4.920) (0.507,5.008) (0.505,5.096) (0.502,5.184) (0.500,5.280)};
  \draw[->] plot[smooth] coordinates {(1.478,2.538) (1.386,2.517) (1.307,2.503) (1.227,2.501) (1.153,2.527) (1.097,2.584) (1.061,2.655) (1.034,2.731) (1.014,2.808) (0.997,2.886) (0.984,2.965) (0.972,3.044) (0.962,3.124) (0.952,3.203) (0.944,3.283) (0.937,3.362) (0.930,3.442) (0.924,3.522) (0.918,3.602) (0.913,3.681) (0.908,3.761) (0.903,3.841) (0.899,3.921) (0.895,4.001) (0.891,4.081)};
  \draw[] plot[smooth] coordinates {(0.891,4.081) (0.887,4.161) (0.884,4.241) (0.880,4.320) (0.877,4.400) (0.874,4.480) (0.871,4.560) (0.869,4.640) (0.866,4.720) (0.863,4.800) (0.861,4.880) (0.859,4.960) (0.856,5.040) (0.854,5.120) (0.852,5.200) (0.850,5.280)};
  \draw[->] plot[smooth] coordinates {(1.478,2.538) (1.389,2.522) (1.327,2.552) (1.294,2.616) (1.273,2.693) (1.260,2.764) (1.249,2.835) (1.241,2.906) (1.232,2.986) (1.226,3.058) (1.220,3.129) (1.215,3.201) (1.210,3.281) (1.206,3.353) (1.202,3.425) (1.198,3.497) (1.195,3.577) (1.192,3.649) (1.189,3.721) (1.186,3.792) (1.183,3.872) (1.181,3.944) (1.178,4.016) (1.176,4.088) (1.174,4.168)};
  \draw[] plot[smooth] coordinates {(1.174,4.168) (1.172,4.240) (1.170,4.312) (1.168,4.384) (1.166,4.464) (1.164,4.536) (1.163,4.608) (1.161,4.680) (1.160,4.760) (1.158,4.832) (1.157,4.904) (1.155,4.976) (1.154,5.056) (1.153,5.128) (1.151,5.200) (1.150,5.280)};
  \draw[->] plot[smooth] coordinates {(1.478,2.538) (1.561,2.564) (1.611,2.615) (1.640,2.681) (1.659,2.750) (1.674,2.821) (1.685,2.892) (1.695,2.963) (1.703,3.035) (1.710,3.106) (1.717,3.178) (1.723,3.250) (1.728,3.321) (1.733,3.393) (1.738,3.465) (1.742,3.537) (1.746,3.609) (1.750,3.681) (1.753,3.753) (1.757,3.825) (1.760,3.905) (1.763,3.977) (1.766,4.048) (1.768,4.120) (1.771,4.192)};
  \draw[] plot[smooth] coordinates {(1.771,4.192) (1.774,4.264) (1.776,4.336) (1.778,4.408) (1.780,4.480) (1.782,4.552) (1.784,4.624) (1.786,4.696) (1.788,4.768) (1.790,4.840) (1.792,4.912) (1.794,4.984) (1.795,5.056) (1.797,5.128) (1.798,5.200) (1.800,5.280)};
  \draw[->] plot[smooth] coordinates {(1.478,2.538) (1.564,2.560) (1.641,2.583) (1.715,2.613) (1.776,2.652) (1.834,2.706) (1.882,2.770) (1.917,2.833) (1.950,2.906) (1.977,2.981) (1.999,3.050) (2.020,3.127) (2.039,3.205) (2.054,3.275) (2.069,3.354) (2.083,3.432) (2.095,3.503) (2.107,3.583) (2.118,3.662) (2.127,3.733) (2.137,3.813) (2.146,3.892) (2.154,3.964) (2.162,4.043) (2.170,4.123)};
  \draw[] plot[smooth] coordinates {(2.170,4.123) (2.176,4.195) (2.183,4.274) (2.190,4.354) (2.196,4.426) (2.202,4.506) (2.208,4.585) (2.213,4.657) (2.218,4.737) (2.223,4.817) (2.228,4.889) (2.233,4.968) (2.237,5.048) (2.241,5.120) (2.246,5.200) (2.250,5.280)};
  \draw[->] plot[smooth] coordinates {(1.478,2.538) (1.576,2.562) (1.661,2.586) (1.745,2.611) (1.829,2.639) (1.903,2.668) (1.984,2.703) (2.062,2.744) (2.136,2.791) (2.200,2.839) (2.266,2.898) (2.325,2.963) (2.379,3.032) (2.427,3.106) (2.467,3.176) (2.506,3.254) (2.541,3.335) (2.573,3.417) (2.599,3.493) (2.625,3.577) (2.650,3.661) (2.672,3.746) (2.691,3.824) (2.710,3.910) (2.728,3.996)};
  \draw[] plot[smooth] coordinates {(2.728,3.996) (2.745,4.082) (2.760,4.169) (2.774,4.248) (2.788,4.335) (2.801,4.422) (2.813,4.509) (2.824,4.588) (2.835,4.675) (2.846,4.763) (2.856,4.850) (2.865,4.930) (2.874,5.017) (2.883,5.105) (2.892,5.192) (2.900,5.280)};
  \draw[->] plot[smooth] coordinates {(1.478,2.538) (1.589,2.566) (1.681,2.591) (1.781,2.621) (1.872,2.650) (1.971,2.683) (2.062,2.715) (2.159,2.750) (2.249,2.784) (2.345,2.823) (2.434,2.861) (2.529,2.904) (2.615,2.946) (2.706,2.995) (2.789,3.044) (2.875,3.102) (2.952,3.160) (3.030,3.228) (3.097,3.297) (3.164,3.376) (3.220,3.454) (3.275,3.542) (3.321,3.627) (3.365,3.721) (3.402,3.810)};
  \draw[] plot[smooth] coordinates {(3.402,3.810) (3.437,3.908) (3.467,3.999) (3.496,4.099) (3.520,4.192) (3.544,4.293) (3.565,4.386) (3.585,4.488) (3.603,4.583) (3.620,4.685) (3.635,4.780) (3.650,4.883) (3.663,4.978) (3.677,5.081) (3.688,5.177) (3.700,5.280)};
  \draw[->] plot[smooth] coordinates {(1.478,2.538) (1.607,2.571) (1.722,2.603) (1.829,2.636) (1.944,2.673) (2.057,2.711) (2.163,2.748) (2.276,2.788) (2.381,2.826) (2.494,2.868) (2.607,2.910) (2.711,2.949) (2.824,2.991) (2.936,3.034) (3.040,3.074) (3.152,3.118) (3.256,3.159) (3.367,3.205) (3.478,3.251) (3.580,3.296) (3.689,3.348) (3.788,3.399) (3.891,3.461) (3.989,3.531) (4.072,3.605)};
  \draw[] plot[smooth] coordinates {(4.072,3.605) (4.150,3.696) (4.215,3.797) (4.265,3.897) (4.308,4.009) (4.340,4.117) (4.369,4.233) (4.393,4.351) (4.413,4.461) (4.431,4.580) (4.445,4.691) (4.459,4.810) (4.471,4.929) (4.481,5.041) (4.491,5.160) (4.500,5.280)};
  \draw[->] plot[smooth] coordinates {(1.478,2.538) (1.391,2.518) (1.313,2.502) (1.234,2.487) (1.155,2.473) (1.076,2.459) (0.998,2.443) (0.921,2.421) (0.848,2.390) (0.782,2.344) (0.729,2.285) (0.686,2.217) (0.653,2.144) (0.626,2.069) (0.603,1.992) (0.583,1.915) (0.566,1.837) (0.550,1.758) (0.536,1.680) (0.523,1.601) (0.512,1.521) (0.500,1.442) (0.490,1.363) (0.480,1.284) (0.470,1.204)};
  \draw[] plot[smooth] coordinates {(0.470,1.204) (0.461,1.125) (0.452,1.045) (0.444,0.966) (0.435,0.886) (0.426,0.807) (0.418,0.727) (0.409,0.647) (0.401,0.568) (0.392,0.488) (0.382,0.409) (0.372,0.330) (0.361,0.250) (0.348,0.172) (0.331,0.093) (0.300,0.020)};
  \draw[->] plot[smooth] coordinates {(1.478,2.538) (1.400,2.520) (1.330,2.504) (1.261,2.482) (1.199,2.446) (1.150,2.394) (1.113,2.332) (1.084,2.266) (1.061,2.198) (1.041,2.128) (1.025,2.058) (1.010,1.988) (0.996,1.917) (0.984,1.846) (0.974,1.783) (0.963,1.712) (0.953,1.641) (0.943,1.569) (0.934,1.498) (0.925,1.426) (0.916,1.355) (0.908,1.283) (0.900,1.212) (0.891,1.140) (0.883,1.069)};
  \draw[] plot[smooth] coordinates {(0.883,1.069) (0.875,0.997) (0.868,0.926) (0.860,0.862) (0.852,0.791) (0.844,0.719) (0.836,0.647) (0.827,0.576) (0.818,0.505) (0.809,0.433) (0.798,0.362) (0.787,0.291) (0.774,0.220) (0.758,0.150) (0.737,0.081) (0.700,0.020)};
  \draw[->] plot[smooth] coordinates {(1.478,2.538) (1.398,2.512) (1.358,2.463) (1.336,2.403) (1.320,2.333) (1.308,2.270) (1.299,2.207) (1.290,2.135) (1.283,2.072) (1.277,2.008) (1.270,1.936) (1.265,1.873) (1.260,1.809) (1.254,1.737) (1.249,1.673) (1.245,1.609) (1.241,1.545) (1.236,1.474) (1.232,1.410) (1.228,1.346) (1.223,1.274) (1.220,1.210) (1.216,1.146) (1.211,1.074) (1.208,1.011)};
  \draw[] plot[smooth] coordinates {(1.208,1.011) (1.204,0.947) (1.199,0.875) (1.195,0.811) (1.191,0.747) (1.187,0.683) (1.182,0.611) (1.177,0.548) (1.173,0.484) (1.167,0.412) (1.161,0.348) (1.155,0.284) (1.147,0.213) (1.138,0.150) (1.126,0.087) (1.100,0.020)};
  \draw[->] plot[smooth] coordinates {(1.478,2.538) (1.561,2.553) (1.622,2.519) (1.651,2.463) (1.672,2.394) (1.687,2.332) (1.700,2.261) (1.712,2.190) (1.721,2.127) (1.730,2.055) (1.738,1.992) (1.746,1.920) (1.754,1.849) (1.761,1.785) (1.768,1.713) (1.774,1.650) (1.780,1.578) (1.787,1.506) (1.793,1.442) (1.799,1.371) (1.805,1.307) (1.811,1.235) (1.816,1.171) (1.822,1.100) (1.829,1.028)};
  \draw[] plot[smooth] coordinates {(1.829,1.028) (1.834,0.964) (1.841,0.892) (1.847,0.829) (1.853,0.757) (1.861,0.685) (1.867,0.622) (1.875,0.550) (1.882,0.487) (1.891,0.415) (1.901,0.344) (1.911,0.281) (1.923,0.210) (1.938,0.147) (1.960,0.079) (2.000,0.020)};
  \draw[->] plot[smooth] coordinates {(1.478,2.538) (1.574,2.562) (1.652,2.580) (1.731,2.592) (1.811,2.586) (1.883,2.553) (1.941,2.499) (1.987,2.433) (2.025,2.363) (2.057,2.289) (2.085,2.214) (2.111,2.139) (2.134,2.062) (2.155,1.985) (2.176,1.908) (2.195,1.830) (2.213,1.752) (2.231,1.674) (2.248,1.596) (2.265,1.518) (2.281,1.439) (2.297,1.361) (2.313,1.283) (2.329,1.204) (2.345,1.126)};
  \draw[] plot[smooth] coordinates {(2.345,1.126) (2.361,1.047) (2.377,0.969) (2.394,0.891) (2.410,0.813) (2.428,0.734) (2.446,0.657) (2.465,0.579) (2.485,0.501) (2.506,0.424) (2.530,0.348) (2.557,0.272) (2.588,0.199) (2.626,0.129) (2.676,0.066) (2.750,0.020)};
  \draw[->] plot[smooth] coordinates {(1.478,2.538) (1.586,2.565) (1.678,2.590) (1.778,2.618) (1.871,2.644) (1.972,2.668) (2.066,2.685) (2.170,2.693) (2.266,2.685) (2.358,2.658) (2.449,2.609) (2.522,2.547) (2.591,2.469) (2.646,2.390) (2.699,2.301) (2.742,2.215) (2.782,2.128) (2.821,2.032) (2.855,1.942) (2.890,1.844) (2.921,1.753) (2.952,1.654) (2.980,1.562) (3.009,1.462) (3.036,1.370)};
  \draw[] plot[smooth] coordinates {(3.036,1.370) (3.062,1.277) (3.090,1.177) (3.116,1.085) (3.143,0.984) (3.169,0.892) (3.198,0.792) (3.225,0.700) (3.253,0.608) (3.286,0.509) (3.318,0.419) (3.356,0.322) (3.396,0.235) (3.449,0.145) (3.511,0.073) (3.600,0.020)};
  \draw[->] plot[smooth] coordinates {(1.478,2.538) (1.619,2.574) (1.742,2.609) (1.872,2.649) (1.994,2.689) (2.115,2.730) (2.244,2.773) (2.366,2.814) (2.487,2.854) (2.617,2.893) (2.741,2.927) (2.866,2.955) (3.000,2.976) (3.128,2.981) (3.255,2.964) (3.375,2.921) (3.487,2.844) (3.574,2.751) (3.650,2.639) (3.710,2.526) (3.761,2.408) (3.808,2.281) (3.848,2.159) (3.884,2.036) (3.918,1.905)};
  \draw[] plot[smooth] coordinates {(3.918,1.905) (3.949,1.780) (3.977,1.655) (4.004,1.530) (4.032,1.397) (4.057,1.272) (4.081,1.146) (4.107,1.012) (4.131,0.887) (4.157,0.753) (4.183,0.628) (4.210,0.503) (4.241,0.370) (4.275,0.247) (4.319,0.127) (4.400,0.020)};
  \draw[->] plot[smooth] coordinates {(1.478,2.538) (1.387,2.517) (1.317,2.503) (1.238,2.488) (1.167,2.476) (1.096,2.466) (1.017,2.455) (0.945,2.446) (0.866,2.438) (0.794,2.431) (0.722,2.425) (0.643,2.419) (0.571,2.415) (0.491,2.411) (0.419,2.408) (0.347,2.405) (0.267,2.402) (0.195,2.400) (0.123,2.392) (0.067,2.342) (0.052,2.271) (0.044,2.192) (0.040,2.120) (0.037,2.048) (0.034,1.968)};
  \draw[] plot[smooth] coordinates {(0.034,1.968) (0.033,1.896) (0.031,1.816) (0.030,1.744) (0.028,1.672) (0.027,1.592) (0.026,1.520) (0.025,1.448) (0.025,1.368) (0.024,1.296) (0.023,1.216) (0.022,1.144) (0.022,1.072) (0.021,0.992) (0.021,0.920) (0.020,0.840)};
  \draw[->] plot[smooth] coordinates {(1.478,2.538) (1.413,2.523) (1.358,2.511) (1.303,2.500) (1.248,2.490) (1.200,2.482) (1.145,2.473) (1.090,2.465) (1.034,2.457) (0.987,2.451) (0.931,2.445) (0.875,2.439) (0.820,2.433) (0.764,2.428) (0.716,2.424) (0.660,2.420) (0.604,2.417) (0.548,2.414) (0.500,2.411) (0.444,2.409) (0.389,2.406) (0.333,2.405) (0.285,2.403) (0.229,2.402) (0.173,2.400)};
  \draw[] plot[smooth] coordinates {(0.173,2.400) (0.117,2.397) (0.065,2.378) (0.046,2.335) (0.038,2.280) (0.033,2.224) (0.030,2.168) (0.028,2.120) (0.027,2.064) (0.025,2.008) (0.024,1.952) (0.023,1.904) (0.022,1.848) (0.021,1.792) (0.021,1.736) (0.020,1.680)};
  \draw[->] plot[smooth] coordinates {(1.478,2.538) (1.640,2.580) (1.794,2.625) (1.946,2.674) (2.098,2.725) (2.241,2.776) (2.392,2.830) (2.542,2.885) (2.692,2.940) (2.834,2.994) (2.984,3.050) (3.134,3.106) (3.284,3.162) (3.434,3.219) (3.576,3.272) (3.726,3.328) (3.875,3.385) (4.025,3.441) (4.167,3.495) (4.317,3.551) (4.467,3.607) (4.617,3.664) (4.759,3.716) (4.910,3.730) (4.946,3.576)};
  \draw[] plot[smooth] coordinates {(4.946,3.576) (4.956,3.416) (4.961,3.256) (4.964,3.104) (4.967,2.944) (4.969,2.784) (4.971,2.624) (4.972,2.472) (4.974,2.312) (4.975,2.152) (4.976,1.992) (4.977,1.840) (4.978,1.680) (4.978,1.520) (4.979,1.360) (4.980,1.200)};
  \draw[->] plot[smooth] coordinates {(1.478,2.538) (1.620,2.574) (1.743,2.610) (1.865,2.647) (1.987,2.687) (2.108,2.729) (2.229,2.771) (2.349,2.814) (2.469,2.858) (2.590,2.902) (2.710,2.947) (2.829,2.992) (2.949,3.037) (3.069,3.082) (3.189,3.127) (3.309,3.172) (3.429,3.217) (3.548,3.262) (3.668,3.307) (3.788,3.352) (3.908,3.397) (4.028,3.442) (4.147,3.487) (4.267,3.532) (4.387,3.577)};
  \draw[] plot[smooth] coordinates {(4.387,3.577) (4.507,3.623) (4.627,3.668) (4.746,3.713) (4.868,3.754) (4.949,3.687) (4.962,3.560) (4.967,3.432) (4.970,3.304) (4.973,3.176) (4.974,3.048) (4.976,2.920) (4.977,2.792) (4.978,2.664) (4.979,2.536) (4.980,2.400)};
  \draw[thick,->] plot[smooth] coordinates {(1.478,2.538) (1.429,2.527) (1.390,2.518) (1.359,2.511) (1.320,2.504) (1.281,2.496) (1.249,2.490) (1.210,2.483) (1.178,2.478) (1.139,2.472) (1.099,2.466) (1.067,2.462) (1.028,2.456) (0.996,2.452) (0.956,2.447) (0.917,2.443) (0.885,2.440) (0.845,2.436) (0.805,2.432) (0.773,2.429) (0.733,2.426)};
  \draw[thick,] plot[smooth] coordinates {(0.733,2.426) (0.701,2.423) (0.662,2.420) (0.622,2.418) (0.590,2.416) (0.550,2.414) (0.518,2.412) (0.478,2.410) (0.438,2.408) (0.406,2.407) (0.366,2.406) (0.326,2.404) (0.294,2.404) (0.254,2.403) (0.222,2.402) (0.182,2.401) (0.142,2.401) (0.110,2.400) (0.070,2.400) (0.030,2.401)};
  \draw[thick,->] plot[smooth] coordinates {(1.478,2.538) (1.579,2.563) (1.671,2.589) (1.763,2.616) (1.855,2.644) (1.946,2.674) (2.037,2.704) (2.128,2.736) (2.211,2.765) (2.301,2.797) (2.392,2.830) (2.482,2.863) (2.572,2.896) (2.662,2.929) (2.752,2.963) (2.842,2.996) (2.924,3.027) (3.014,3.061) (3.104,3.095) (3.194,3.128) (3.284,3.162)};
  \draw[thick,] plot[smooth] coordinates {(3.284,3.162) (3.374,3.196) (3.463,3.230) (3.553,3.264) (3.636,3.295) (3.726,3.328) (3.815,3.362) (3.905,3.396) (3.995,3.430) (4.085,3.464) (4.175,3.497) (4.265,3.531) (4.347,3.562) (4.437,3.596) (4.527,3.630) (4.616,3.664) (4.706,3.698) (4.796,3.732) (4.886,3.766) (4.970,3.810)};
  % fixed points
  \fill[blue] (0,0) circle (0.07); \fill[blue] (\sx,0) circle (0.07);
  \fill[green!60!black] (0,\sy) circle (0.07); \fill[green!60!black] (0.2956*\sx,0) circle (0.07);
  \fill[green!60!black] (\sx,1.5874*\sy) circle (0.07);
  \fill[red] (0.2956*\sx,1.0574*\sy) circle (0.08);
  \node[left,font=\small] at (0,0) {$0$}; \node[left,font=\small] at (0,\sy) {$1$};
  \node[left,font=\small] at (-0.3,1.6*\sy) {$y$};
  \node[below,font=\small] at (0,0) {$0$}; \node[below,font=\small] at (0.5*\sx,0) {$0.5$}; \node[below,font=\small] at (\sx,0) {$1$};
  \node[below=10pt,font=\small] at (0.65*\sx,0) {$z$};
  \node[blue,font=\small,anchor=north east] at (-0.05,-0.1) {$(0,0)$};
  \node[blue,font=\small,anchor=north west] at (\sx+0.05,-0.1) {$(1,0)$};
  \node[green!60!black,font=\small,anchor=east] at (-0.35,\sy) {$(0,1)$};
  \node[green!60!black,font=\small,anchor=north] at (0.2956*\sx,-0.35) {$(0.30,0)$};
  \node[green!60!black,font=\small,anchor=west] at (\sx+0.05,1.5874*\sy) {$(1,1.59)$};
  \node[red,font=\small,anchor=west] at (\sx+0.05,1.0574*\sy) {$(0.30,1.06)$};
\end{tikzpicture}
\diagramcaption{Schematic RG flow of the spin-1 Ising model in the plane $x = 0$. Color coding: blue,
attractive along both eigenvectors (can be 1-sided); green, attractive along one eigenvector and repulsive
along the other eigenvector; red, repulsive along both eigenvectors. Here $d\nu = 2-\alpha$,
$d\nu' = 2-\alpha'$.}
\end{diagram}

From the RG transformation, we can characterize each of these RGFPs. To do this, we linearize the RG
transformation at each fixed point and calculate the eigenvectors and eigenvalues. Write
\begin{equation}
  z' = R_1(z,y) = \frac{4z^2}{(1+z^2)^2},\qquad y' = R_2(z,y) = \frac{y^4}{(1+z^2)^2}.
\end{equation}
Then
\begin{equation}
  M = \begin{pmatrix}\dfrac{\partial z'}{\partial z} & \dfrac{\partial z'}{\partial y}\\[6pt] \dfrac{\partial y'}{\partial z} & \dfrac{\partial y'}{\partial y}\end{pmatrix}
    = \begin{pmatrix}\dfrac{8z(1-z^2)}{(1+z^2)^3} & 0\\[10pt] \dfrac{-4zy^4}{(1+z^2)^3} & \dfrac{4y^3}{(1+z^2)^2}\end{pmatrix}.
\end{equation}
If an eigenvalue $>1$, the direction is repulsive; if an eigenvalue $<1$, attractive; if an eigenvalue $=1$,
marginal.

\paragraph{RGFP at the lower left corner, $(z_*,y_*) = (0,0)$:} $M = \left(\begin{smallmatrix}0&0\\0&0\end{smallmatrix}\right)$;
we can choose $\hat{\mathbf e}_1 = \hat{\mathbf e}_z = \binom{-1}{0}$, $\lambda_z = 0$, attractive;
$\hat{\mathbf e}_2 = -\hat{\mathbf e}_y = \binom{0}{-1}$, $\lambda_y = 0$, attractive.

\paragraph{RGFP at the lower right corner, $(z_*,y_*) = (1,0)$:}
$M = \left(\begin{smallmatrix}0&0\\0&0\end{smallmatrix}\right)$, so $\lambda_1 = 0$, $\lambda_2 = 0$: attractive
along both eigenvectors, $\hat{\mathbf e}_1 = \binom10$, $\hat{\mathbf e}_2 = \binom{0}{\pm1}$. So the RGFPs
$(0,0)$ and $(1,0)$ are stable to perturbations along both $\hat{\mathbf e}_1$ and $\hat{\mathbf e}_2$.

The eigenvalues of a triangular matrix are the diagonal entries. So we have
\begin{equation}
  \lambda_2 = \frac{4y^3}{(1+z^2)^2},\quad \hat{\mathbf e}_2 = \binom01\ \text{(vertical direction)};\qquad
  \lambda_1 = \frac{8z(1-z^2)}{(1+z^2)^3};
\end{equation}
for $M = \left(\begin{smallmatrix}\lambda_1&0\\a&\lambda_2\end{smallmatrix}\right)$,
$\hat{\mathbf e}_1\propto\binom{a^{-1}(\lambda_1-\lambda_2)}{1}$.

\paragraph{RGFP at $(z_*,y_*) = (0,1)$:} $M = \left(\begin{smallmatrix}0&0\\0&4\end{smallmatrix}\right)$
$\Rightarrow\lambda_1 = 0$, $\hat{\mathbf e}_1 = \binom10$: attractive along the horizontal direction (i.e.\
stable); $\lambda_2 = 4 = 2^{y_2}$, $\hat{\mathbf e}_2 = \binom01$: repulsive along the vertical direction (i.e.\
unstable), so $y_2 = 2$ (don't confuse the exponents $y_i$ with the Boltzmann variable $y$). Thus, in the local
vicinity of this RGFP, if there is a slight deviation along the vector $\hat{\mathbf e}_1$, i.e.\
$(z,y) = (\epsilon,1)$, then under the iterated RG mapping the point flows back toward the RGFP at $(0,1)$, but
if the initial point is $(0,1\pm\delta)$, i.e.\ above or below $(0,1)$, then the RG maps this point further
away, either upward or downward, in a vertical direction.

\paragraph{RGFP at the bottom ``middle'', $(z_*,y_*) = (0.30,0)$:}
$M = \left(\begin{smallmatrix}1.679&0\\0&0\end{smallmatrix}\right)$; $\lambda_1 = 1.679$,
$\hat{\mathbf e}_1 = \binom10$; $\lambda_1>1$, so repulsive in the horizontal direction $\hat{\mathbf e}_1$;
$\lambda_2 = 0$, $\hat{\mathbf e}_2 = \binom01$; $\lambda_2<1$, so attractive in the vertical direction. With
$\lambda_1 = b^{y_1}\Rightarrow\ln\lambda_1 = y_1\ln b$,
\begin{equation}
  y_1 = \frac{\ln\lambda_1}{\ln b} = \frac{\ln(1.679)}{\ln 2} = 0.7473 .
\end{equation}

\paragraph{RGFP at the right ``middle'', $(z_*,y_*) = (1,2^{2/3}) = (1,1.587)$:}
$\Rightarrow M = \left(\begin{smallmatrix}0&0\\-2^{5/3}&2^2\end{smallmatrix}\right)$; $\lambda_1 = 0$,
attractive along $\hat{\mathbf e}_1$, an oblique direction as shown in the diagram; $\lambda_2 = 2^2$,
repulsive along $\hat{\mathbf e}_2 = \binom01$, the vertical direction; $\lambda_2 = 2^{y_2}\Rightarrow y_2 = 2$.

\paragraph{RGFP in the interior, $(z^*,y^*) = (0.2956,1.057)$:} this can be analyzed in a similar manner.
